\section{Rubric-Gated Oracle-Free Testing of Mesh-Based Neural
Surrogates for Cylinder
Flow}\label{rubric-gated-oracle-free-testing-of-mesh-based-neural-surrogates-for-cylinder-flow}

\begin{quote}
Manuscript draft for IST regular track.\\
Draft status: v0.3 evidence-gated full draft scaffold, 2026-06-05.\\
Results remain blocked until executable MR artifacts and experiment
ledgers exist.
\end{quote}

\subsection{Target and Scope}\label{target-and-scope}

\begin{itemize}
\tightlist
\item
  Primary target: \emph{Information and Software Technology} regular
  research-paper track.
\item
  Paper type: empirical software engineering / software testing and V\&V
  method paper.
\item
  Review framing: software testing contribution first; SciML and
  cylinder flow are the application context.
\item
  Word limit target: keep the submitted paper below 15,000 words.
\item
  Current draft policy: write Background, Motivation, Method, Empirical
  Design, and Threats; do not write empirical findings before artifacts
  exist.
\end{itemize}

\subsection{Structured Abstract}\label{structured-abstract}

\subsubsection{Context}\label{context}

Scientific machine-learning (SciML) surrogates are increasingly used to
approximate expensive physical simulations. Mesh-based neural simulators
are attractive for fluid-flow problems because they operate on irregular
meshes and predict spatiotemporal physical fields. Testing such systems
is difficult because exact expected outputs are often unavailable for
arbitrary inputs without running a trusted high-fidelity solver. Rollout
accuracy and physics residuals provide useful evidence, but they do not
by themselves specify which physical relations should hold under
controlled transformations of inputs, meshes, boundary conditions, or
parameters.

\subsubsection{Objective}\label{objective}

This paper investigates how physically meaningful metamorphic relations
(MRs) can be identified, screened, and operationalized as executable
oracle-free tests for MeshGraphNets-family cylinder-flow surrogates. The
goal is not to improve the predictive model, but to make relation-based
validation more systematic, auditable, and explicit about the physical
and numerical conditions under which each relation is expected to hold.

\subsubsection{Method}\label{method}

We propose a domain-validity rubric for screening candidate MRs, a
NOETHER-informed workflow for organizing candidate relation structures,
and an executable MR asset format that records source cases, follow-up
transformations, output mappings, metrics, tolerances, exclusion rules,
and relation-level verdicts. We design an empirical evaluation over
three MeshGraphNets-family implementations or configurations, comparing
the proposed workflow with expert-designed MRs, generic MR-generation
scope contrasts, LLM-assisted candidate generation, and rollout-accuracy
baselines.

\subsubsection{Results}\label{results}

BLOCKED. Empirical findings will be drafted only after executable MR
cards, SUT configurations, run logs, and experiment ledgers are
available. This draft makes no claims about pass rates, fault-detection
rates, comparative performance, localization accuracy, or SUT
reliability.

\subsubsection{Conclusion}\label{conclusion}

At the planning stage, the intended contribution is a validity-aware
testing workflow for SciML surrogates. The paper positions MRs as a
complement to rollout accuracy, residuals, uncertainty estimates, and
equivariance errors: these diagnostics can become executable oracle-free
relations only when paired with physically valid transformations,
explicit preconditions, and auditable verdict rules.

\subsection{Keywords}\label{keywords}

Metamorphic testing; metamorphic relation identification; oracle
problem; scientific machine learning; MeshGraphNets; cylinder flow;
software verification and validation.

\subsection{1. Introduction}\label{introduction}

Scientific computing increasingly uses learned surrogates to reduce the
cost of repeated numerical simulation. In fluid dynamics, mesh-based
neural simulators such as MeshGraphNets learn to predict physical fields
on irregular meshes and can provide fast rollout predictions for
benchmark problems such as flow around a cylinder. These models are
useful because they can approximate high-cost solvers, but they also
create a familiar software testing problem in a new setting: for many
candidate inputs, there is no cheap and exact expected output against
which the program can be checked.

The common response is to evaluate a surrogate on held-out trajectories
and report rollout error. This is necessary, but it is not the whole
validation problem. A model may have acceptable average error on a
finite test set while still failing to preserve relations that should
hold under a physically meaningful transformation. Conversely, a high
error value does not always explain which physical, numerical, or
representational condition has been violated. For a user who wants to
deploy a SciML surrogate outside a narrow validation set, the practical
question is not only ``How accurate is the model on these samples?'' but
also ``Under which transformations or regimes does this system stop
respecting the relations it should respect?''

Metamorphic testing (MT) addresses the test-oracle problem by checking
relations among multiple executions instead of requiring an exact output
for each individual test case. A metamorphic relation states that if a
source input is transformed in a specified way, the corresponding
outputs should satisfy a necessary relation. This is a natural idea for
scientific computing: conservation laws, symmetry relations,
nondimensional similarity, boundary constraints, continuity conditions,
and numerical stability expectations all suggest possible relations
among executions.

The difficult part is deciding which candidate relations are valid and
executable for a particular SciML program. A transformation that looks
plausible at a high level may violate the governing assumptions,
boundary conditions, discretization choices, or measurement tolerance of
the actual system under test (SUT). For cylinder flow, mirror symmetry
depends on geometry and boundary labels; translation invariance depends
on how coordinates and domain boundaries are represented;
Reynolds-Strouhal similarity depends on nondimensionalization and flow
regime; a divergence check depends on a discrete divergence operator,
mesh weights, and boundary treatment. Treating such transformations as
automatically valid MRs would make the test suite look rigorous while
hiding the very assumptions that determine whether a violation is
meaningful.

Recent work has started to connect physics-based metamorphic testing
with learned physical-field predictors. The closest identified studies
use physical correlations, disturbance scenarios, cross-domain
scenarios, Reynolds-number conservation, and Strouhal-number
correlations to test physical-field or fluid-velocity prediction models.
These papers show that physics-based MR ideas are already useful for
learned physical systems. They also sharpen the gap addressed here. The
remaining problem is not whether physical knowledge can inspire MRs; it
is how candidate MRs are screened for domain validity, turned into
executable test assets, and reported through relation-level verdicts
that distinguish SUT inconsistency from out-of-relation-domain cases and
numerical tolerance effects.

This paper treats MR identification for SciML as a validity-gated
testing problem. We use the NOETHER framework only as an organizing
structure for candidate relation patterns. It does not prove the
physical validity of a candidate MR. Physical validity is established,
or rejected, through an explicit rubric that records the relation's
physical basis, boundary-condition compatibility, transformation
preconditions, output mapping, metric, tolerance rationale, and
exclusion rule. Only retained relations are converted into executable MR
assets.

The study is scoped to MeshGraphNets-family cylinder-flow surrogates.
This is an intentionally narrow scope. It allows us to examine
physically meaningful transformations over meshes, geometry, velocity
fields, nondimensional quantities, and rollout behavior while keeping
the SUT family concrete enough for reproducible testing. The evaluation
is planned across three implementations or configurations, but we do not
claim external validity across all neural fluid surrogates.

\subsubsection{1.1 Research Questions}\label{research-questions}

The main research question is:

\textbf{RQ0. How can physically valid metamorphic relations be
identified, operationalized, and used as oracle-free tests for
MeshGraphNets-family cylinder-flow surrogates without relying on exact
per-sample expected outputs?}

We decompose this into four questions.

\textbf{RQ1 - Validity.} How can a domain-validity rubric distinguish
physically meaningful candidate MRs from transformations that are
executable but invalid, underspecified, or outside the relation's
domain?

\textbf{RQ2 - Operationalization.} How can a NOETHER-informed candidate
organization workflow be combined with the rubric to produce executable
MR assets with source cases, follow-up cases, metrics, thresholds, and
verdict rules?

\textbf{RQ3 - Verdict and interpretation.} How can relation-level
verdicts distinguish pass, fail, skip, out-of-relation-domain, numerical
tolerance issue, and inconclusive outcomes?

\textbf{RQ4 - Empirical comparison.} Across three MeshGraphNets-family
cylinder-flow implementations or configurations, how do retained MRs
compare with expert-designed MRs, generic MR-generation scope contrasts,
LLM-assisted candidate generation, and rollout-accuracy baselines?

\subsubsection{1.2 Contributions}\label{contributions}

This paper is planned around four contributions.

First, we propose a domain-validity rubric for screening candidate MRs
in SciML testing. The rubric checks whether a candidate relation has a
clear physical basis, compatible boundary conditions, a
semantics-preserving transformation, a measurable output relation, an
interpretable tolerance, and a diagnosable failure mode.

Second, we present a NOETHER-informed and rubric-gated workflow for
candidate MR organization and operationalization. NOETHER is used as a
candidate-generation and organization framework, while physical validity
is decided by the rubric, execution constraints, and evidence records.

Third, we define an executable MR asset format and relation-level
verdict scheme. Each retained MR records a source-case schema, follow-up
transformation, output mapping, metric, tolerance rule, exclusion rule,
and verdict interpretation. This converts physical diagnostics such as
residuals, conservation errors, and equivariance errors into auditable
oracle-free tests only when their transformation and validity conditions
are explicit.

Fourth, we design an empirical evaluation over three
MeshGraphNets-family cylinder-flow implementations or configurations.
The evaluation will compare retained MRs with expert MR design, generic
MR-generation scope contrasts, LLM-assisted candidate generation, and
rollout-accuracy baselines. This contribution will be stated as a result
only after experiment artifacts exist.

\subsection{2. Background and Related
Work}\label{background-and-related-work}

\subsubsection{2.1 Mesh-Based Neural
Simulation}\label{mesh-based-neural-simulation}

Mesh-based neural simulators learn dynamics on graph or mesh
representations of physical systems. For fluid-flow problems, a mesh
representation is attractive because the geometry, boundary regions, and
local connectivity can be represented without forcing the state onto a
regular grid. MeshGraphNets is a representative architecture family in
this area: it encodes mesh nodes and edges, propagates information by
message passing, and predicts future physical fields through
autoregressive rollout.

The cylinder-flow benchmark is useful for testing because it combines
several features that matter for SciML validation. It involves geometry,
boundary conditions, velocity fields, pressure-like quantities or
derived quantities, temporal rollout, and flow-regime assumptions. It
also exposes the distinction between data-driven accuracy and physical
relation preservation. A model may fit a trajectory distribution while
failing under a controlled transformation of geometry, mesh
representation, or flow parameters.

In this paper, MeshGraphNets-family systems are treated as software
systems under test. We do not evaluate them as new modeling
contributions. Instead, we ask how their outputs behave under controlled
input transformations and whether necessary relations remain within
explicit tolerances.

\subsubsection{2.2 Metamorphic Testing and the Oracle
Problem}\label{metamorphic-testing-and-the-oracle-problem}

Metamorphic testing was developed to address programs for which it is
difficult or impossible to know the correct output for a single test
input. Instead of checking one output against one expected value, MT
checks a relation among the outputs of a source input and one or more
follow-up inputs. This framing has been applied to scientific software,
simulation models, and machine-learning systems, all of which can suffer
from oracle problems.

For SciML surrogates, the oracle problem is especially visible in
out-of-distribution (OOD) settings. A trusted solver may be too
expensive to run for every transformed case, and even when a reference
output is available, a single pointwise error does not necessarily
indicate whether a physical relation has been preserved. MR-based
testing offers a complementary perspective: it asks whether the SUT
maintains a necessary relation under a controlled transformation.

However, SciML MRs cannot be treated as generic input perturbations. In
image-based ML testing, transformations such as lighting or weather
changes may preserve the semantic label. In scientific computing, the
transformation must respect governing equations, physical regimes,
boundary conditions, mesh representation, and numerical tolerance. The
correctness of the MR itself becomes a central validity question.

\subsubsection{2.3 MR Identification for Scientific and Simulation
Software}\label{mr-identification-for-scientific-and-simulation-software}

Prior work on scientific software testing has shown that MRs can be
elicited from monotonicity, conservation-like behavior, scaling
relations, simulation assumptions, and domain-specific expectations.
Work on simulation validation similarly treats multi-run relations as
pseudo-oracles when direct validation data are limited or unavailable.
These studies establish that MR thinking is practical for scientific and
simulation software, not merely for toy functions.

Other work has explored automated or semi-automated MR identification.
For example, research on ocean system models shows that data-driven
search can discover candidate transformations in complex physical
software. CPS testing work shows that design assumptions can be encoded
as MRs and then used to generate new test traces. These ideas are
important for the present paper because they show that MR sources can
include system assumptions, model structure, and physical design
knowledge.

The gap is that candidate identification is not enough for SciML. A
candidate relation must also state when it is supposed to hold. Without
a domain-of-validity record, a failed test may mean several different
things: the relation was invalid under the chosen boundary conditions,
the transformation left the physical regime, the tolerance was not
numerically justified, or the SUT violated a relation it should have
preserved. This paper makes that distinction explicit.

\subsubsection{2.4 Physics-Based MT for Learned Scientific
Simulators}\label{physics-based-mt-for-learned-scientific-simulators}

Recent studies are directly relevant because they test learned
physical-field or fluid-velocity predictors using physics-based
metamorphic ideas. In particular, work on multi-scenario testing for
physical-field prediction models and work on mutation testing for
fluid-velocity prediction models use physical correlations and
dimensionless quantities such as Reynolds and Strouhal numbers to
construct follow-up scenarios and evaluate physical consistency or
extrapolation.

These studies are the closest neighbors of this paper. They show that
physics-derived relations can guide testing of learned physical
predictors, including perturbation and cross-domain scenarios. This
paper therefore does not claim that physics-based MT for learned fluid
or field predictors is new.

Our contribution is narrower and methodological. We shift the emphasis
from scenario-level physical consistency checks to a validity-gated
workflow. A retained relation is represented as an executable MR asset,
not only as an intuitive physical expectation. Each asset records its
preconditions, boundary-condition compatibility, output mapping, metric,
tolerance rationale, exclusion rule, and verdict interpretation. This
allows the test report to distinguish a model inconsistency from a
relation that was applied outside its valid domain.

\subsubsection{2.5 SciML V\&V, Residuals, UQ, and Failure
Modes}\label{sciml-vv-residuals-uq-and-failure-modes}

SciML reliability research has developed several important tools:
physics residuals, uncertainty quantification, conformal prediction,
certified error bounds, stress testing, equivariance error, and
failure-mode analysis. These tools address real weaknesses of learned
PDE solvers and neural surrogates. For example, calibrated
physics-informed uncertainty quantification uses PDE residuals as
nonconformity scores and provides statistical coverage guarantees. Work
on PINN failure modes shows that physical constraints in a training loss
do not guarantee reliable behavior in more difficult regimes.

These methods are complementary to MT. A residual, conservation error,
or equivariance error can be a useful diagnostic, but it is not
automatically an MR. It becomes part of an executable relation only when
there is a source case, a follow-up transformation, an expected output
relation, a tolerance rule, and a verdict interpretation. This
distinction is central to our paper. We use SciML diagnostics as
possible relation measurements, while MT supplies the multi-execution
oracle-free structure.

\subsubsection{2.6 Hybrid ML-Solver Trust
Regions}\label{hybrid-ml-solver-trust-regions}

Hybrid ML-solver frameworks provide another relevant line of work. Some
systems use residual thresholds or error estimates to decide when a
learned component should be trusted and when a numerical solver should
take over. Such systems operationalize a form of trust region, although
often at runtime and through residual-based switching rather than
offline relation-based testing.

Our study pursues a complementary direction. Before deployment,
physically derived transformations are used to estimate where relation
violations occur and what regimes, boundary conditions, or numerical
assumptions are associated with them. The result is not a runtime
switching policy by itself, but an evidence structure that can support
later decisions about when a learned surrogate should be trusted.

\subsection{3. Method}\label{method-1}

\subsubsection{3.1 Overview}\label{overview}

The proposed method is a five-stage workflow:

\begin{enumerate}
\def\labelenumi{\arabic{enumi}.}
\tightlist
\item
  identify candidate relation sources;
\item
  organize candidate relations using a NOETHER-informed structure;
\item
  screen candidates with a domain-validity rubric;
\item
  convert retained relations into executable MR assets;
\item
  execute the assets and record relation-level verdicts.
\end{enumerate}

The method deliberately separates candidate generation from validity
judgment. NOETHER may help organize or generate candidate relation
structures, but it does not certify that a relation is physically valid
for a particular SUT, dataset, mesh, boundary condition, or numerical
tolerance. Validity is determined by the rubric and by executable
evidence.

\subsubsection{3.2 Candidate Relation
Sources}\label{candidate-relation-sources}

For cylinder-flow surrogates, candidate MRs may come from six sources.

\textbf{Physical equations and constraints.} Incompressible continuity,
boundary behavior, and conservation-like expectations can suggest
relation measurements such as discrete divergence or boundary-condition
consistency.

\textbf{Geometric and symmetry assumptions.} Mirroring, rigid
transformations, and coordinate-frame changes may suggest equivariance
or invariance checks, but only under compatible geometry and boundary
labels.

\textbf{Nondimensional similarity.} Reynolds-number and Strouhal-number
relations may suggest follow-up transformations over flow parameters or
extracted wake behavior, but only within regimes where the empirical or
theoretical relation is meaningful.

\textbf{Mesh and graph representation.} Node permutations, face
ordering, edge encoding, and mesh refinement may suggest
representation-level MRs.

\textbf{Temporal rollout behavior.} Autoregressive rollouts may suggest
determinism, prefix consistency, or semigroup-like sanity checks. These
are useful implementation or numerical-consistency checks, but should
not be overstated as physical laws.

\textbf{Cross-implementation comparison.} Outputs from different
implementations may support method-comparison checks only if units,
state variables, meshes, rollout horizons, boundary conditions, and
checkpoints are comparable. Otherwise they are triangulation evidence
rather than retained MRs.

\subsubsection{3.3 Domain-Validity Rubric}\label{domain-validity-rubric}

Each candidate MR is screened using a rubric with six criteria.

\textbf{Physical basis.} The relation must cite its source: governing
equation, boundary condition, nondimensional law, representation
property, numerical assumption, or implementation contract.

\textbf{Transformation preconditions.} The source-to-follow-up
transformation must state what is changed and what is preserved. For
example, a mirror relation must specify how coordinates, vector
components, node types, and boundary labels transform.

\textbf{Boundary-condition compatibility.} The relation must hold under
the boundary conditions of the transformed case. If the transformation
changes the physical meaning of a boundary, the candidate is rejected or
marked as an OOD stress test rather than a physics-preserving MR.

\textbf{Output mapping.} The expected relation among outputs must be
measurable. A relation that cannot be mapped to available output fields
is not executable.

\textbf{Metric and tolerance.} The verdict rule must define a metric and
a tolerance. The tolerance may come from numerical precision, solver
reference behavior, calibration data, repeated-run variability, or
expert thresholding, but the source must be recorded.

\textbf{Failure diagnosability.} The relation should indicate what kind
of failure or boundary condition it can help interpret. If every
possible violation has the same ambiguous meaning, the relation is weak
evidence and should be treated cautiously.

The rubric outputs one of four screening decisions:

\begin{itemize}
\tightlist
\item
  retained as executable MR;
\item
  retained as OOD stress relation, not physics-preserving MR;
\item
  rejected as invalid or underspecified;
\item
  deferred pending missing evidence, such as a discrete operator or
  threshold.
\end{itemize}

\subsubsection{3.4 MR Card and Executable Asset
Format}\label{mr-card-and-executable-asset-format}

A retained MR is represented as an MR card. The card records:

\begin{itemize}
\tightlist
\item
  MR identifier and name;
\item
  source category;
\item
  physical or software basis;
\item
  source-case schema;
\item
  follow-up transformation;
\item
  transformation preconditions;
\item
  boundary-condition compatibility;
\item
  output mapping;
\item
  metric;
\item
  tolerance and provenance;
\item
  expected verdict classes;
\item
  exclusion rules;
\item
  artifact schema;
\item
  expected fault or boundary interpretation.
\end{itemize}

Executable assets are generated from MR cards. Each asset includes
transformation code, runner configuration, metric computation, verdict
logic, and artifact recording. The asset format is intended to make MR
execution auditable: another researcher should be able to inspect why a
relation was retained, how the follow-up case was generated, how the
comparison was made, and what a violation means.

\subsubsection{3.5 Relation-Level
Verdicts}\label{relation-level-verdicts}

An MR execution can produce several verdicts:

\begin{itemize}
\tightlist
\item
  \textbf{pass:} the relation holds within the stated tolerance;
\item
  \textbf{fail:} the relation is violated within its valid domain;
\item
  \textbf{skip:} the case cannot be run because a declared precondition
  is missing;
\item
  \textbf{out-of-relation-domain:} the transformation produced a case
  outside the relation's validity domain;
\item
  \textbf{numerical-tolerance issue:} the evidence is dominated by
  threshold or numerical-resolution uncertainty;
\item
  \textbf{inconclusive:} the artifact is insufficient to decide.
\end{itemize}

This scheme prevents a common overclaim: not every relation violation is
a program fault. A violation may reveal a model inconsistency, but it
may also reveal that the MR was applied outside its domain, that the
threshold was too tight, that the mesh operator was unsuitable, or that
the case is not comparable.

\subsubsection{3.6 Hierarchical Interpretation
Protocol}\label{hierarchical-interpretation-protocol}

We organize retained MRs into a hierarchy inspired by
scientific-computing MR classification: physical-model relations,
computational-model relations, and code-model relations. For learned
mesh surrogates, the computational-model level includes graph
representation and message-passing discretization assumptions.

The hierarchy supports a predeclared interpretation protocol.
Representation-level MRs may point to graph encoding or adapter
problems. Physical-model MRs may point to continuity, symmetry, or
similarity violations. Code-model MRs may point to determinism, rollout,
or implementation issues. At this stage, this is a localization
protocol, not a validated localization model. It becomes validated only
if seeded faults or mutants with known layers are used to evaluate the
inference rule.

\subsection{4. Empirical Design}\label{empirical-design}

\subsubsection{4.1 Subject Systems}\label{subject-systems}

The planned study uses three MeshGraphNets-family implementations or
configurations:

\begin{enumerate}
\def\labelenumi{\arabic{enumi}.}
\tightlist
\item
  an echowve MeshGraphNets PyTorch/PyG implementation;
\item
  NVIDIA PhysicsNeMo \texttt{vortex\_shedding\_mgn};
\item
  DeepMind TF1 MeshGraphNets, or an equivalent third
  MeshGraphNets-family implementation or configuration if runtime
  feasibility requires substitution.
\end{enumerate}

For each SUT, the experiment ledger will record repository URL, commit,
checkpoint, dataset, mesh format, input fields, output fields, rollout
horizon, random seeds, environment, and known runtime limitations.
Because these SUTs share a task and architecture family, the study will
be framed as a same-family stress test, not as evidence for all neural
fluid surrogates.

\subsubsection{4.2 Planned MR Classes}\label{planned-mr-classes}

The initial MR set will be grouped into six classes:

\textbf{Representation MRs:} node permutation, face-order invariance,
and encoding consistency.

\textbf{Geometric and symmetry MRs:} mirror-y equivariance and rigid
transformation candidates, retained only under boundary-compatible
conditions.

\textbf{Continuity and constraint MRs:} discrete divergence or
mass-continuity checks with explicit mesh weighting and boundary
treatment.

\textbf{Nondimensional similarity MRs:} Reynolds-Strouhal or scaling
candidates, retained only under nondimensional and regime-compatibility
checks.

\textbf{Numerical stability MRs:} small perturbation stability, mesh
perturbation, or refinement consistency candidates.

\textbf{Temporal and implementation MRs:} rollout prefix consistency,
deterministic execution, and conditional cross-implementation
comparison.

Time reversal is excluded as a retained MR for viscous cylinder flow.
Divergence is treated as an incompressible continuity or
mass-conservation constraint, not as a Noether-derived conservation law.

\subsubsection{4.3 Baselines and
Comparators}\label{baselines-and-comparators}

The study uses four comparator families.

\textbf{Expert MR design.} Domain experts manually propose MRs. This
baseline estimates what the proposed rubric adds beyond conventional
expert elicitation.

\textbf{Generic MR-generation scope contrast.} Generic MR identification
or generation methods are applied as far as their assumptions allow.
This comparator is not used to claim that generic methods are weak; it
is used to show where domain-validity information is needed for SciML.

\textbf{LLM-assisted candidate generation.} An LLM generates candidate
MRs from prompts containing equations, SUT descriptions, and candidate
relation categories. The LLM is treated only as a candidate-generation
and material-organization tool. It does not judge MR validity. Prompt
logs, model/version, temperature, candidate lists, and rubric decisions
will be recorded.

\textbf{Rollout-accuracy baseline.} Standard rollout error is reported
to compare pointwise predictive evidence with relation-level evidence.
The goal is not to prove that MRs are better than accuracy, but to
identify whether they answer a different validation question.

Where feasible, residual-based or UQ metrics will be reported as
diagnostic comparators. They will not be treated as MRs unless paired
with a source/follow-up transformation and explicit verdict rule.

\subsubsection{4.4 Efficacy Parameters and
Metrics}\label{efficacy-parameters-and-metrics}

The primary efficacy parameters are:

\begin{itemize}
\tightlist
\item
  candidate retention rate: the proportion of candidate MRs retained by
  the rubric;
\item
  executable MR rate: the proportion of retained MRs that can be
  executed on each SUT;
\item
  MR violation rate: the proportion of valid-domain executions that fail
  the relation;
\item
  violation magnitude: the metric distance from the stated tolerance;
\item
  fault-detection rate: the proportion of seeded faults or mutants
  detected by at least one retained MR;
\item
  boundary characterization: the transformation regions or parameter
  bins associated with elevated violation rates;
\item
  interpretation yield: the proportion of failed executions assigned a
  non-ambiguous verdict category.
\end{itemize}

Secondary parameters include:

\begin{itemize}
\tightlist
\item
  MR construction cost;
\item
  inter-rater agreement for rubric decisions;
\item
  flakiness rate across repeated executions;
\item
  localization agreement with seeded-fault layers;
\item
  complementarity with rollout accuracy, measured by cases where
  accuracy and MR verdicts provide different boundary information.
\end{itemize}

\subsubsection{4.5 Statistical Plan}\label{statistical-plan}

Violation rates will be reported with confidence intervals. Paired
comparisons will be used where the same source cases and transformations
are run across SUTs or baselines. For binary verdict comparisons,
McNemar or Fisher-style tests may be used when assumptions are
appropriate. For continuous violation magnitudes, paired bootstrap
intervals are preferred. When many MRs or transformation bins are
compared, multiple-comparison correction or false-discovery-rate control
will be reported.

If the number of executable cases is small, the study will emphasize
effect sizes, confidence intervals, and qualitative failure
interpretation rather than strong significance claims.

\subsubsection{4.6 Ethics, Integrity, and
Reproducibility}\label{ethics-integrity-and-reproducibility}

The planned study does not involve human subjects, personal data, or
private sensitive information. The main ethical and integrity issues are
research transparency, AI use, and reproducibility.

All SUT versions, datasets, checkpoints, scripts, and run logs should be
recorded when licensing permits. Failed, skipped, and inconclusive cases
must remain in the experiment ledger. LLM use is restricted to candidate
generation and evidence organization; it is not used as a final judge of
MR validity. No candidate MR should be described as valid unless it
passes the rubric and has a corresponding MR card. No OOD violation
should be described as a program fault without the verdict evidence
needed to support that interpretation.

Third-party code, datasets, and model checkpoints will be used according
to their licenses. If any artifact cannot be redistributed, the paper
will provide access instructions and record the limitation.

\subsection{5. Results}\label{results-1}

BLOCKED until experiments exist.

This section must not contain:

\begin{itemize}
\tightlist
\item
  pass/fail rates;
\item
  comparative superiority claims;
\item
  fault-detection rates;
\item
  localization accuracy;
\item
  runtime or performance claims;
\item
  claims that one SUT is more reliable than another.
\end{itemize}

Permitted placeholders:

\begin{itemize}
\tightlist
\item
  table shells for MR cards;
\item
  planned result tables;
\item
  artifact requirements;
\item
  analysis templates;
\item
  expected interpretation rules.
\end{itemize}

\subsection{6. Discussion}\label{discussion}

\subsubsection{6.1 Expected
Interpretation}\label{expected-interpretation}

The expected value of the study is not that every MR will find a new
fault beyond rollout accuracy. Rather, the value is that retained MRs
provide relation-level evidence under explicit transformations. If a
violation is observed, the MR card and verdict rule should help
determine whether the issue is a model inconsistency, a relation-domain
boundary, a numerical tolerance problem, or an inconclusive case.

\subsubsection{6.2 Boundary of Claims}\label{boundary-of-claims}

The paper must avoid four overclaims. It must not claim that MRs are
always better than rollout accuracy. It must not claim that NOETHER
proves the physical validity of cylinder-flow MRs. It must not claim
that LLMs can automatically identify valid MRs. It must not generalize
from three MeshGraphNets-family configurations to all SciML surrogates.

The safer claim is that a domain-validity-aware workflow can make MR
identification and execution more auditable for a concrete class of
SciML SUTs.

\subsubsection{6.3 Implications for SciML
Testing}\label{implications-for-sciml-testing}

If the empirical study succeeds, it will show how SciML testing can move
from implicit expert checks to explicit MR assets. Such assets can
complement residuals, uncertainty estimates, and accuracy metrics by
making transformation assumptions and verdict rules inspectable. This is
especially important for OOD validation, where the boundary of the
relation is often as important as the violation itself.

\subsection{7. Threats to Validity}\label{threats-to-validity}

\textbf{Construct validity.} MR validity depends on the rubric, physical
assumptions, boundary-condition compatibility, and tolerance rules.
Incorrect discrete operators or thresholds may create false violations
or false passes.

\textbf{Internal validity.} SUT setup, checkpoint differences, random
seeds, mesh preprocessing, and runtime nondeterminism may affect
verdicts. The experiment ledger must record these details.

\textbf{External validity.} The study is limited to MeshGraphNets-family
cylinder-flow implementations or configurations. It should not be
generalized to all neural operators, PINNs, or fluid surrogates without
further evidence.

\textbf{Baseline fairness.} Generic MR-generation and LLM baselines may
not be designed for SciML. They should be interpreted as scope contrasts
and candidate-generation comparators, not as defeated competitors.

\textbf{Conclusion validity.} Small sample sizes, multiple MRs, and many
transformation bins can produce unstable conclusions. The analysis
should emphasize confidence intervals, effect sizes, and predeclared
verdict categories.

\textbf{Reproducibility.} Some SUTs may depend on old runtimes or
non-redistributable checkpoints. The paper should disclose such barriers
and provide the most complete runnable package possible.

\subsection{8. Conclusion}\label{conclusion-1}

This draft does not provide a final conclusion because empirical
evidence is not yet available. The intended conclusion, if supported by
artifacts, is that domain-validity-gated MR identification can provide
an auditable oracle-free testing workflow for MeshGraphNets-family
cylinder-flow surrogates. The central claim will remain methodological:
physically meaningful SciML MRs require explicit validity conditions,
executable assets, and relation-level verdicts.

\subsection{References}\label{references}

Reference verification is ongoing. Planned reference groups:

\begin{itemize}
\tightlist
\item
  foundational MT and the oracle problem;
\item
  MT for ML and scientific software;
\item
  MT for physical-field or fluid-velocity predictors;
\item
  SciML V\&V, residuals, UQ, and failure modes;
\item
  MeshGraphNets and mesh-based neural simulation;
\item
  NOETHER and the scientific-computing MR classification model.
\end{itemize}
